% GENERATED FILE. Edit canonical lessons, then run npm run generate:books.
% canonical-source-hash: fnv1a64:d427c8a130422e03
% canonical-lessons: JA-W145-four-lines

\chapter{Writing: Four Lines, and the Joint You Do Not Write}
\label{ch:ja-writing-four-lines}
\hlchaptermodality{145}

\begin{chapteropening}
\textbf{By the end of this chapter:} \emph{I can write four connected lines of Japanese to a named friend, walking through my day with time words and naming who does it at most once.}

Complete JA-W145-four-lines: a greeting and three lines of your day in order, each opening on a time word and closing on a polite verb, then count how many lines name who does it.
\end{chapteropening}

\section[(four lines)]{Four lines, and the joint you do not write}

\label{lesson:JA-W145-four-lines}



\begin{quote}
The last lesson asked for writing \textbf{finished on time}. This one asks for writing that \textbf{holds together}.

Pick one \textbf{\ja{ともだち}} who will read it.
\end{quote}

\subsection*{What to know first}
Connected writing is not \textbf{longer} writing. Four short lines can be connected; one long line can fail to be. What makes it connected is that a line stops standing on its own.

In many languages the joints are little words between the sentences. This book has not taught you those for Japanese, and you do not need them yet. Your four lines are held together by two things you already own.

\textbf{The first is a time word at the front of each line.} \textbf{\ja{あさ}}, then \textbf{\ja{ゆうがた}}: the reader follows your day because each line says where in it they are.

\textbf{The second is the joint you do not write.} Chapter 138 taught you that Japanese leaves out who does it when that is clear. Say who once, or not at all, and every line after it leans on the one before. Name yourself in every line and you have written a \textbf{list}, not a message.

\subsection*{Writing — connected composition}
Write \textbf{four lines} to the \textbf{\ja{ともだち}} you picked. The first greets them with \textbf{\ja{こんにちは}}. The other three walk through your day in order, each opening on a time word and closing on a polite verb: \textbf{\ja{おきます}}, \textbf{\ja{はたらきます}}, \textbf{\ja{かえります}} or \textbf{\ja{ねます}}.

Use only words you already write. This lesson adds no vocabulary; the whole exercise is the joining.

End every line with \textbf{\ja{。}}, a small open circle, not a dot.

\begin{tcolorbox}[breakable,colback=teal!4,colframe=teal!35!black,title={Before you move on}]
Read your third line alone. It names no one: who comes home, or goes to bed? Only the lines before it say. It leans on its neighbours, and that is what connected means.

Then count how many lines name who does it. If every line does, you wrote a list, which is where you were one lesson ago.
\end{tcolorbox}
